% ============================================================
\chapter{Modul 14: Machine Learning Grundlagen}
% ============================================================

\section{Lektion 38: Was ist Machine Learning?}

\begin{analogie}
Klassische Programmierung: Du beschreibst die Regeln. Machine Learning: Du zeigst Beispiele, das System lernt die Regeln selbst. Statt \texttt{if farbe == "rot": stop} zeigst du tausend Bilder und sagst: ``Das hier war gefährlich, das nicht.''
\end{analogie}

In der Robotik: Objekterkennung, Anomalie-Erkennung in Sensorwerten, Lokalisierung, Bahnplanung.

\section{Lektion 39: Erstes Modell -- Hindernisklassifikation}

\begin{lstlisting}
import numpy as np
from sklearn.preprocessing import StandardScaler
from sklearn.neighbors import KNeighborsClassifier
from sklearn.model_selection import train_test_split
from sklearn.metrics import accuracy_score

# Trainingsdaten: [abstand_vorne, abstand_links, abstand_rechts]
# Label: 0=frei, 1=links_ausweichen, 2=rechts_ausweichen, 3=stopp
np.random.seed(42)
n = 500

X = np.random.uniform(0.1, 3.0, (n, 3))
y = np.zeros(n, dtype=int)

for i, (vorne, links, rechts) in enumerate(X):
    if vorne < 0.3:
        y[i] = 3   # stopp
    elif vorne < 0.6 and rechts > links:
        y[i] = 1   # links ausweichen
    elif vorne < 0.6 and links >= rechts:
        y[i] = 2   # rechts ausweichen

# Trainieren
X_train, X_test, y_train, y_test = train_test_split(
    X, y, test_size=0.2, random_state=42)

scaler = StandardScaler()
X_train_s = scaler.fit_transform(X_train)
X_test_s  = scaler.transform(X_test)

clf = KNeighborsClassifier(n_neighbors=5)
clf.fit(X_train_s, y_train)

# Bewerten
y_pred = clf.predict(X_test_s)
print(f"Genauigkeit: {accuracy_score(y_test, y_pred)*100:.1f}%")

# Echtzeit-Vorhersage
aktuelle_sensoren = np.array([[0.45, 0.8, 0.3]])   # Hindernis vorne
aktuelle_s = scaler.transform(aktuelle_sensoren)
aktion = clf.predict(aktuelle_s)[0]
aktionen = {0: "Frei fahren", 1: "Links ausweichen",
            2: "Rechts ausweichen", 3: "STOPP"}
print(f"Empfohlene Aktion: {aktionen[aktion]}")
\end{lstlisting}

\section{Lektion 40: Vorhersagen im Echtzeitsystem}

\begin{lstlisting}
import joblib   # pip install joblib

# Modell speichern (nach dem Training)
joblib.dump(clf,    "hindernis_clf.pkl")
joblib.dump(scaler, "scaler.pkl")

# Modell laden (im Roboterprogramm)
clf_geladen    = joblib.load("hindernis_clf.pkl")
scaler_geladen = joblib.load("scaler.pkl")

class ML_Hindernisvermeidung:
    def __init__(self, clf_pfad, scaler_pfad):
        self.clf    = joblib.load(clf_pfad)
        self.scaler = joblib.load(scaler_pfad)

    def entscheide(self, abstand_vorne, abstand_links, abstand_rechts):
        X = np.array([[abstand_vorne, abstand_links, abstand_rechts]])
        X_s = self.scaler.transform(X)
        return self.clf.predict(X_s)[0]
\end{lstlisting}

\begin{merksatz}
Trainieren ist offline (langsam, viel Rechenzeit). Vorhersagen ist online (schnell, läuft auf dem Roboter). Das Modell wird einmal trainiert, dann auf den Roboter übertragen und dort geladen.
\end{merksatz}

\begin{uebung}[title=Projektaufgabe Modul 14: Gesten-Erkennung]
Trainiere einen Klassifikator für 4 Fahrzustände (vorwärts, rückwärts, links, rechts) basierend auf 5 Sensorwerten. Generiere synthetische Trainingsdaten, trainiere, bewerte mit Accuracy und Confusion Matrix (aus \texttt{sklearn.metrics}).
\end{uebung}
